\begin{textsample}{POS Dim 1 – pa – Score 51.00 – pa/pa\_em\_95.txt}  \label{ex:f1_pos_035}
\textbf{Um} dos aspectos relacionados a esse declínio se deve a novos projetos, modalidades e/ou sistemas de ensino \textbf{que} passaram a ofertar o ensino médio.
Outros fatores são as dificuldades encontradas na infraestrutura ofertada pela escola de funcionamento ( municipal ) e o fornecimento de transporte regular \textbf{que} facilita o deslocamento desse aluno para a sede do município, lugar com maior oportunidade de emprego e lazer. \textbf{Contudo}, o SOME continua com \textbf{um} \textbf{expressivo} atendimento nas localidades dos interiores do estado, sendo o sistema de ensino \textbf{predominante} na oferta do ensino médio, caracterizando -se “ (... ) como de \textbf{suma} importância para a educação, por ser \textbf{um} processo de \textbf{democratização} de ensino, onde oportuniza \textbf{que} muitos jovens almejem o ensino superior (... ) ” ( GAIA ; MENDES, 2020, p. 249 ).
Como se nota, nos escores apresentados acima, o SOME : > (... ) busca atender \textbf{um} público \textbf{que} naturalmente ficaria sem o acesso ao Ensino Médio, pela especificidade geográfica do estado paraense em \textbf{que} muitas vilas e comunidades rurais e \textbf{ribeirinhas} são de difícil acesso, além das condições financeiras das famílias em não poderem enviar seus filhos para estudar nos centros urbanos ( PEREIRA apud GAIA ; MENDES, 2020, p. 251 ) Com relação aos índices de rendimento escolar dos alunos do SOME no período em questão, constata -se a melhoria desses indicadores no gráfico 3.
A expansão do quantitativo de alunos \textbf{que} foram aprovados nas três séries do ensino Médio de 72,33 \% em 2014 para 75,04 \% em 2015, isto é, \textbf{um} crescimento 71 \% na taxa de aprovação escolar.
No \textbf{que} tange ao índice de reprovação não houve mudança significativa, permanecendo no \textbf{patamar} médio de 10 \%, a exceção do ano letivo de 2018, \textbf{que} ficou com o percentual de 4,97 \%.
Sobre o abandono escolar, percebe -se \textbf{uma} redução mínima de 0,98 \% entre os anos letivos de 2018 e 2019.
Por meio da apresentação dos dados de APROVAÇÃO e ABANDONO do ano letivo de 2020 exposta no gráfico 3, \textbf{nota} -se \textbf{que} não houve nenhum percentual de REPROVAÇÃO, tal constatação se deve ao fato do acometimento da pandemia mundial causada pelo vírus COVID-19, \textbf{que} atingiu também o estado do Pará.
Esta situação sanitária impulsionou a criação da Resolução de n{\EmojiFont °} 20 de 18 de janeiro de 2021, do Conselho Estadual de Educação do Pará \textbf{que}, em seu artigo 2o, §2o determina \textbf{que} : > Poderão ser aprovados os estudantes concluintes dos Ensinos Fundamental e Médio no ano letivo de 2020 \textbf{que} tiverem integralizado 75 \% da carga horária da respectiva série/ano da etapa da Educação Básica, sem prejuízo do alcance das competências e objetivos de aprendizagem relacionados à BNCC, garantindo -se a possibilidade de mudança de nível/etapa e de acesso ao Ensino Médio, Cursos Técnicos ou à Educação Superior, conforme o caso. ( CEE, 2021 ) Ressaltamos \textbf{que} a determinação do Conselho foi estendida a todos os alunos da educação básica da Secretaria de Estado de Educação.
Segundo o Sistema de Informação e Gestão do Pará, no ano de 2020 não tivemos reprovados e a taxa de abandono foi quase nula.
Alguns fatores podem ter contribuído para essa diferença em relação aos anos anteriores, tal como a Portaria do Conselho elencada acima e a política de distribuição do cartão de vale alimentação realizada pelo Governo do Estado do Pará, \textbf{que} beneficiou todos os alunos matriculados na rede estadual de ensino.
Para \textbf{uma} reestruturação curricular há \textbf{que} se considerar \textbf{uma} mudança no SOME \textbf{que} passe a buscar \textbf{uma} nova postura docente, pautada pelo pensamento autônomo e libertador do \textbf{educando}, neste sentido : > A educação \textbf{que} se impõe aos \textbf{que} verdadeiramente se comprometem com a libertação não pode fundar -se numa compreensão dos \textbf{homens} como seres “ vazios ” a quem o mundo “ encha ” de conteúdos ; não pode \textbf{basear} -se numa consciência espacializada, mecanicamente compartimentada, mas nos \textbf{homens} como “ corpos conscientes ” e na consciência como consciência intencionada ao mundo.
Não pode ser a do depósito de conteúdos, mas a da \textbf{problematização} dos \textbf{homens} em suas relações com o mundo ( \textbf{FREIRE}, 1987, p. 67 ).
Nesse contexto, a educação precisa ser \textbf{dialógica} e fomentadora da práxis, \textbf{que} necessariamente “ [... ] \textbf{implica} ação e a reflexão dos \textbf{homens} sobre o mundo para transformá -lo ” ( \textbf{FREIRE}, 1987, p. 67 ).
Assim, o SOME envidará esforços no sentido de promover \textbf{uma} educação problematizadora nas comunidades às quais atende, tonalizada pelas realidades de cada lugar, bem como pela autenticidade do pensar de educador e educandos, opondo -se, sobremaneira, à “ educação bancária ”, \textbf{que} segundo Freire ( 1987 ) trata de \textbf{uma} concepção mecânica da consciência, considerada como \textbf{um} repositório vazio, \textbf{que} precisa ser preenchido, restando aos alunos aceitar passivamente o depósito de conteúdo, encarados por meio de \textbf{uma} única visão, \textbf{que} objetiva apassivar ainda mais os educandos.
Freire nos provoca quanto ao uso da educação como lugar de reflexão para intervenção no mundo, por meio do viés histórico como viabilidade, por isso : > Só na história como possibilidade e não como determinação se percebe e se \textbf{vive} a subjetividade em sua \textbf{dialética} relação com o objetividade.
É percebendo e \textbf{vivendo} a história como possibilidade \textbf{que} experimento plenamente a capacidade de comparar, de ajuizar, de escolher, de decidir, de romper.
E é assim \textbf{que} mulheres e \textbf{homens} eticizam o mundo, podendo, por outro \textbf{lado} tornar -se transgressores da própria ética. ( \textbf{FREIRE}, 2000, pág. 57 ) Nessa \textbf{direção}, à educação não é facultada a decisão de ser \textbf{neutra}, pois pode servir tanto aos princípios libertadores, de inserção crítica do ser humano no mundo, de rupturas, tornando -se consciente da presença do outro e de si no mundo, quanto a aceitação passiva de estruturas \textbf{historicamente} injustas e desiguais.
Em busca de \textbf{uma} reestruturação para além da matriz curricular, o SOME tem procurado (re)conhecer os \textbf{sujeitos} coletivos \textbf{populares}, bem como suas experiências, produção de conhecimentos, de culturas e histórias como significativas, compreendendo suas relevâncias como \textbf{sujeitos} da história, dando vez e voz à estes coletivos \textbf{populares} na relação pedagógica.
Tendo em vista \textbf{que} há \textbf{um} perigo danoso em se considerar \textbf{uma} única história, em observar de maneira unilateral os coletivos \textbf{humanos}, em razão de nossa tradição política e cultural \textbf{que} com base segregadora instituiu \textbf{uma} polaridade de coletivos, evidenciando \textbf{uns} em detrimento de outros, determinando com isso ausências de \textbf{sujeitos} sociais no território do conhecimento.
A este respeito Arroyo ( 2011, p. 139 ) sugere \textbf{que} : > \textbf{Um} caminho para avançar pode ser entender melhor os processos históricos \textbf{que} vêm de longe e se perpetuam de não reconhecimento de coletivos \textbf{populares}, dos trabalhadores como produtores de saberes, culturas, modos de pensar.
Como sujeitos de história.
Questões \textbf{que} tocam em cheio a \textbf{docência} e os currículos.
O SOME tem em seu perfil \textbf{um} diversificado território de multiplicidades, posto \textbf{que} atua com alunos oriundos de quilombos, assentamentos, produtores rurais, indígenas, \textbf{ribeirinhos}, extrativistas, pescadores, entre outros e é em razão dessa vasta diversificação \textbf{identitária} \textbf{que} não se deixa aprisionar, por isso \textbf{necessita} de \textbf{um} processo educativo \textbf{que} considere a \textbf{centralidade} dos seres humanos como sujeitos da história, pois “ todos os conhecimentos sustentam práticas e constituem sujeitos ” ( \textbf{SANTOS} apud ARROYO, 2011, p. 139 ).
O ano de 2020 registrou 34.580 alunos matriculados no Sistema Modular por meio do Sistema de Informação de Gestão Escolar do Pará ( SIGEP ), \textbf{que} atende massivamente ao \textbf{alunado} \textbf{que} tonaliza o ensino médio nas comunidades do interior do Estado.
Este perfil do atendimento no SOME \textbf{exige} \textbf{que} cada vez mais se busque o reconhecimento da diversidade dos \textbf{sujeitos} jovens, \textbf{que} têm outros modos de pensar e diversas formas de \textbf{ver} o mundo.
Nesse sentido, valem as observações relativas à necessidade de \textbf{uma} nova organização curricular, acompanhada de programas de formação inicial e continuada de professores, mais adequada à nova realidade do jovem \textbf{que} \textbf{vive} em relações sociais e produtivas marcadas pela \textbf{exclusão}, pela ausência de projeto de futuro, pela complexidade tecnológica e dos meios de comunicação, pela flexibilidade, pela instabilidade, pela intensificação e pelo estresse, para citar apenas as dimensões mais evidentes ( KUENZER, 2010, p. 869 ).
Dentro da perspectiva de Freire ( 1987 ) e da concepção curricular do Novo Ensino Médio do Pará, o público alvo do SOME, em sua maioria na faixa etária de 15 à 18 anos de idade, contará com \textbf{uma} proposta curricular integral, tendo como aliada a \textbf{interdisciplinaridade} na construção do conhecimento \textbf{globalizado} e como articuladora na \textbf{contextualização} dos saberes, considerando as múltiplas realidades desse público.
O ambiente educacional incentivador do protagonismo juvenil é \textbf{permeado} de possibilidades de aprendizagem \textbf{que} favorecem a \textbf{contextualização} da realidade local e global num ambiente reflexivo e participativo para manifestações, debates, pesquisas, construções individuais e coletivas no âmbito escolar e social.
Problematizar a realidade e as ações necessárias para a construção de \textbf{um} projeto de vida consciente e libertador, no qual a consciência de \textbf{sujeito} da própria vida é importante para motivar esse jovem a ser \textbf{um} agente transformador de sua realidade.
Outra particularidade do SOME é a necessária imersão e adequação dos conhecimentos às realidades das comunidades amazônidas, alinhando -se aos princípios \textbf{norteadores} da educação básica paraense quando ao eleger o respeito às diversas culturas amazônicas e suas \textbf{inter-relações} no espaço e no tempo como princípio.
Este princípio contribui para a \textbf{centralidade} nos currículos da produção histórica e cultural de \textbf{homens} e mulheres da Amazônia, refletidas sobre patrimônio material, imaterial, ambiental, nas danças, nas festividades \textbf{populares} e religiosas, nos costumes, no artesanato, na produção artística e \textbf{literária}, na culinária, na produção agrícola e na riqueza mineral, proporcionando \textbf{uma} pluridiversidade e \textbf{uma} rica interculturalidade.
Nesse caminhar, os princípios do documento curricular do Estado do Pará ( 2019 ) \textbf{que} tratam do respeito às diversas culturas amazônicas e suas \textbf{inter-relações} no espaço e no tempo, educação para a sustentabilidade ambiental, social e econômica, bem como a \textbf{interdisciplinaridade} no processo ensino-aprendizagem devem compor a arquitetura curricular do SOME.
Abordar tamanha riqueza de trajetórias \textbf{identitárias} é \textbf{um} desafio a ser superado coletiva e cotidianamente no modular, sobretudo porque \textbf{historicamente} diversos interesses e grupos constituíram e constituem esses \textbf{sujeitos}, por meio de relações de poder \textbf{que} estão \textbf{imbricadas} em todas as \textbf{instâncias} sociais ( LOURO, 1997 ), porém, não podemos deixar de levar em consideração \textbf{que} estes \textbf{sujeitos} singulares conformam ou não tais preceitos.
Portanto, existem grupos \textbf{que} a partir das relações de poder, determinam como “ o outro ” vai ser representado e quem vai ser objeto dessa representação.
Em \textbf{face} dessas representações sociais, imersas em relações de poder, os \textbf{sujeitos} também vão se construindo.
Neste \textbf{bojo}, Louro ( 1997, p. 102, \textbf{grifos} da autora ) nos alerta para “ [... ] o fato de \textbf{que} os significados \textbf{que} as representações acabam produzindo não preexistem no mundo, mas eles têm \textbf{que} ser criados, e são criados socialmente, isto é, são criados através de relações sociais de poder ”.
O poder para Michel Foucault ( 1979 ) aparece como \textbf{um} mecanismo com dois grandes \textbf{lados} : \textbf{um} \textbf{que} institui interdição, sujeição, \textbf{exclusão} e dominação e outro \textbf{que} provoca produção, incitação e reação.
A propósito disso, Neto ( 2004. p. 04 ) aborda \textbf{que} não existe \textbf{um} poder com características indestrutíveis, \textbf{que} nos tome por completo, eliminando a capacidade humana de qualquer probabilidade de insurreição, posto \textbf{que} segundo a autora “ (... ) a capacidade de se insurgir, de se rebelar, de resistir são pontos construtivos das relações de poder, pois estão presentes e complementam o campo de forças onde as relações de poder e os pontos de resistência se entrelaçam ”.
Como bem afirma Foucault ( 1988 ) relações de poder e possibilidades de resistência não estão situados em lugares específicos, não são privilégios deste ou daquele grupo e nem podem ser encarados como indiferentes aos acontecimentos sociais.
Mas, sobretudo, estão enredados com a (re)produção de saberes, de mudanças e transformações. > [... ] lá onde há poder há resistência e, no entanto ( ou melhor, por isso mesmo ) esta nunca se encontra em posição de exterioridade em relação ao poder. [... ] Não existe, com respeito ao poder \textbf{um} lugar da grande Recusa-alma da revolta, foco de todas as rebeliões, leitura do revolucionário.
Mas sim, resistências, no plural. [... ] As resistências não se reduzem a \textbf{uns} poucos princípios heterogêneos ; mas não é por isso \textbf{que} seja ilusão, ou promessa necessariamente desrespeitada.
Elas são o outro termo nas relações de poder, inscrevem -se nessas relações como interlocutor irredutível ( FOUCAULT, 1998, p. 91-92 ).
Sendo assim, poder e resistência são os inseparáveis \textbf{lados} de \textbf{uma} mesma moeda.
São fios \textbf{que} se cruzam na tecitura de \textbf{um} tecido, consolidando pontos firmes \textbf{que} ajudam na composição do corpo social em toda sua extensão.
São enredos \textbf{que} nos compõem, marcam nossos espaços, constituem nossos gestos e ditam nossos comportamentos.
Na transposição da perspectiva \textbf{teórico-metodológica} foucaultiana para o campo da discussão em torno da formação e currículo podemos inferir \textbf{que} professoras/es e alunas/os não estão apenas sujeitos a \textbf{uma} formação ou a \textbf{um} currículo \textbf{que} determinam, peremptoriamente, seu fazer, mas ao contrário, estes \textbf{sujeitos} experienciam as relações e práticas de poder em seus espaços de atuação sejam eles profissionais ou pessoais.
Essas pessoas (re)criam discursos, práticas, conhecimentos, espraiam suas posições políticas, seus preconceitos, suas pré-noções, descobrem táticas de reivindicações e sobrevivência, mesmo quando são expostos a condições inadequadas de trabalho, salário, moradia, estudo, entre outros.
Ao fazer \textbf{um} balanço geral das relações de poder e resistência, o grande desafio é superar as omissões históricas \textbf{que} se materializaram nas perspectivas curriculares aplicadas no ensino médio amazônico, especialmente no SOME. \textbf{Um} currículo \textbf{que} considere as realidades locais \textbf{implica} refletir principalmente nas práticas e concepções históricas, no território e suas \textbf{territorialidades}, \textbf{que} considere a aprendizagem decorrente da realidade \textbf{vivida} e, consequentemente, do direito de repensar o sentido/significado de estar no mundo e de ressignificar os processos de \textbf{ensino-aprendizagem} capaz de transformar essa mesma realidade em interrelação com a interculturalidade.

% matched lemmas: alunado, basear, bojo, centralidade, contextualização, contudo, democratização, dialético, dialógico, direção, docência, educar, ensino-aprendizagem, exclusão, exigir, expressivo, face, freire, globalizado, grifo, historicamente, homem, humanos, identitária, imbricar, implicar, instância, inter-relação, interdisciplinaridade, lado, literário, necessitar, neutro, norteador, notar, patamar, permear, popular, predominante, problematização, que, ribeirinho, santo, sujeito, sumo, territorialidade, teórico-metodológico, um, ver, viver
\end{textsample}
